\section*{Asesmen}
\addcontentsline{toc}{section}{Asesmen}

Asesmen Bab XIV menjadi bagian Praktikum 4 (8\%) dan Tugas Teori
(8\%). Fokusnya adalah SDD, translation scheme, dan inferensi tipe.

\begin{table}[H]
\centering
\caption{Instrumen asesmen Bab XIV}
\label{tab:ch14-asesmen}
\begin{tabularx}{\textwidth}{@{}lXl@{}}
\toprule
\textbf{Instrumen} & \textbf{Indikator} & \textbf{Bobot} \\
\midrule
P4-A & Inferensi tipe ekspresi & bagian P4 \\
Tugas Teori & SDD dan translation scheme & 8\% \\
Kuis & Atribut synthesized/inherited & formatif \\
\bottomrule
\end{tabularx}
\end{table}

Kriteria kelulusan: mahasiswa dapat menulis SDD, menjelaskan
perbedaan synthesized dan inherited, serta mengimplementasikan
inferensi tipe dasar.
